% Я шаблонный код вынес в другой файл и импортировал его сюда
\documentclass[a4paper, 12pt]{article}

% золотой набор преамблы
\usepackage[english, russian]{babel} % Подключает языковые модули. Настраивает правила переноса слов и автоматически переводит системные слова (например, вместо "Contents" пишет "Оглавление", вместо "Theorem" — "Теорема"). Последний язык в списке (russian) становится основным.

\usepackage[T2A]{fontenc} % Задает внутреннюю кодировку шрифтов LaTeX. T2A — это стандартная кодировка для поддержки кириллицы. Благодаря ей в итоговом PDF-файле русский текст будет нормально копироваться и искаться.

\usepackage[utf8]{inputenc} % Говорит компилятору, что твой исходный код .tex написан в кодировке UTF-8. Как я говорил ранее, в современном Overleaf этот пакет уже встроен по умолчанию, его можно опускать.

\usepackage{amsmath, amsfonts, amssymb, amsthm, mathtools} % интерфейс работы с теоремами обеспечивается пакетом amsthm
\usepackage{graphicx} % это главный пакет в LaTeX для вставки и редактирования изображений.
\usepackage[left=2cm, right=2cm, top=2cm, bottom=2cm]{geometry}
\usepackage{hyperref} % работа с ссылками, делает кликабельным
\hypersetup{
    colorlinks=true,    % Включает цветной текст ссылок вместо уродливых рамок
    linkcolor=blue,     % Цвет внутренних ссылок (оглавление, формулы)
    urlcolor=cyan,      % Цвет внешних веб-ссылок
    citecolor=green,    % Цвет ссылок на литературу
    pdftitle={Конспект по матанализу}, % Метаданные для PDF-файла
    pdfauthor={Студент 1-го курса}
}

\usepackage{indentfirst} % пакет для абзаца каждой секции
\usepackage{titleps} % пакент для колонтитулов
\usepackage{soulutf8} % пакет для разного вида шрифта
\usepackage{xcolor} % для работы с цветами

% для работы с таблицами
\usepackage{array}
\usepackage{tabularx}
\usepackage{multirow}

% стиль страницы
\newpagestyle{main}{ % main - название стиля, и в других фигурных скобках это содержимое
    \setheadrule{0.4pt} % толщина линии отделяющий хедер или футер от текста
    \sethead{ФПМИ Конспект}{}{Глава \thesection} % верхний колонтитул (header) обязательные три параметра!
    % \thesection — выведет номер раздела (например, 1.2)
    % \sectiontitle — автоматически выведет имя текущего раздела
    \setfootrule{0.4pt}  % толщина линии отделяющий хедер или футер от текста
    \setfoot{Конспект лекций}{\thepage}{2026} % нижний колонтитул (footer)
}
\pagestyle{main} % команда для приминения стиля страницы, можно использовать в каком-то определенном промежутке кода

% стиль plain
\theoremstyle{plain} % модификатор, который задает стиль всем следуюшим теоремам, пока не переопределить.
\newtheorem{theorem}{Теорема}
\newtheorem{lemma}{Лемма}
\newtheorem{axeom}{Аксиома}
\newtheorem{corollary}{Следствие}[theorem] % последний необязательный аргумент, говорит о нумерации следствия данной теоремы.

% стиль defenition
\theoremstyle{definition}
\newtheorem{definition}{Определение}[theorem] % привязанность к данной теореме
\newtheorem{property}{Свойство}[theorem]
\newtheorem{exercise}{Задача}[section]

% стиль remark
\theoremstyle{remark}
\newtheorem{remark}[theorem]{Замечание} % здесь спрятана подготовительная информация

\begin{document}
    Hello World! % здесь пишется сам текст
% Привет Мир!
    
    Equations type of: $y = x$
    $$ax = b$$
% \newline
    $$y = ax^2$$ % это центрированная формула, то есть по центру строки
    \[ax^2+ bx + c = 0\] % это центрированная формула, то есть по центру строки, другой записи
    Now, math formulas will be shown: $\sqrt{2} + \sqrt{2} = 2\sqrt{2}$\\ % квадратный корень
    Optional argument of root: $\sqrt[3]{2}$\\
    Экранированный символ, то есть просто фигурные или любые другие скобки: $\{ \}$\\
    Команда просто точка и команда крестик: $\cdot   \times$\\
    Команда рисующая число пи: $\pi$\\
    Integral: \[\int_{-\infty}^{+\infty} e^{\frac{x^2}{2}} dx = \sqrt{2 \pi}\]
    
    Индексы: \[x_n, x^n, x_i^n, x_{i+j}^{n+k}, x_{n_k}\]
    
    Греческий алфавит: \[\alpha, \beta, \gamma, \delta, \epsilon, \phi, \Delta\]
    \[\forall, \exists\] и так далее
    
    Бинарные операции: \[\times, \ge, \le, =, \lt, \gt, \cdot, \cong \oplus, \sim \]
    \[\frac{\frac{1}{2} + \frac{3}{4}}{\frac{1}{8}} = \frac{5}{4}\cdot10\]
    
    Операторы: $sin x$, $\mathrm{sin}{x}$, $\sin{x}$, $\mathrm{cos}{x}$
    
    Скобки: \[(\frac{\pi}{2})\], $\big(\frac{\pi}{2}\big)$, $\bigg(\frac{\pi}{2}\bigg)$, $\left(\frac{\pi}{2}\right)$, $\left\{\frac{\pi}{2}\right\}$, $\left\[\frac{\pi}{2}\right]$
    
    Если скобка нужна только одна, то после left или right пишется просто точка:
    \[
        \left\{
            \begin{aligned} % окружение aligned или можно matrix Чтобы уравнения внутри системы стояли аккуратно друг под другом, их принято заворачивать в дополнительное мини-окружение оно работает как раз благодаря пакету amsmath
                &\sin{x} = 1 \\ % \\ переносит на новую строку, а & выравнивает уравнения он ставится перед началом каждого уравнения внутри системы, чтобы выровнять их строго по левому краю (по одной вертикальной линии).
                &\cos{x} = 0
            \end{aligned}
        \right.
    \]
    
    \[ % Открываем математику
        \begin{aligned}
            x + y &= 10 \\  % Выравнивание пойдет строго по знаку &
            2x    &= 5
        \end{aligned}
    \] % Закрываем математику
    
    $\mathrm{\sin}{\left(\frac{\pi}{2}\right)}$
    
    Как экранировать символы:
    \backslash,
    \$,
    \{,\},
    \&,
    \#,
    \_,
    \%
\end{document}